%document
\documentclass[10pt]{beamer}
%theme
\usetheme{CambridgeUS}
% packages
\usepackage{color}
\usepackage{listings}
\usepackage[ngerman]{babel}
\usepackage[utf8]{inputenc}
\usepackage{multicol}


% color definitions
\definecolor{mygreen}{rgb}{0,0.6,0}
\definecolor{mygray}{rgb}{0.5,0.5,0.5}
\definecolor{mymauve}{rgb}{0.58,0,0.82}

\lstset{
    backgroundcolor=\color{white},
    % choose the background color;
    % you must add \usepackage{color} or \usepackage{xcolor}
    basicstyle=\footnotesize\ttfamily,
    % the size of the fonts that are used for the code
    breakatwhitespace=false,
    % sets if automatic breaks should only happen at whitespace
    breaklines=true,                 % sets automatic line breaking
    captionpos=b,                    % sets the caption-position to bottom
    commentstyle=\color{mygreen},    % comment style
    % deletekeywords={...},
    % if you want to delete keywords from the given language
    extendedchars=true,
    % lets you use non-ASCII characters;
    % for 8-bits encodings only, does not work with UTF-8
    frame=single,                    % adds a frame around the code
    keepspaces=true,
    % keeps spaces in text,
    % useful for keeping indentation of code
    % (possibly needs columns=flexible)
    keywordstyle=\color{blue},       % keyword style
    % morekeywords={*,...},
    % if you want to add more keywords to the set
    numbers=left,
    % where to put the line-numbers; possible values are (none, left, right)
    numbersep=5pt,
    % how far the line-numbers are from the code
    numberstyle=\tiny\color{mygray},
    % the style that is used for the line-numbers
    rulecolor=\color{black},
    % if not set, the frame-color may be changed on line-breaks
    % within not-black text (e.g. comments (green here))
    stepnumber=1,
    % the step between two line-numbers.
    % If it's 1, each line will be numbered
    stringstyle=\color{mymauve},     % string literal style
    tabsize=4,                       % sets default tabsize to 4 spaces
    % show the filename of files included with \lstinputlisting;
    % also try caption instead of title
    language = Java,
	showspaces = false,
	showtabs = false,
	showstringspaces = false,
	escapechar = ,
}

\def\ContinueLineNumber{\lstset{firstnumber=last}}
\def\StartLineAt#1{\lstset{firstnumber=#1}}
\let\numberLineAt\StartLineAt



\newcommand{\codeline}[1]{
	\alert{\texttt{#1}}
}

% This Document contains the information about this course.

% Authors of the slides
\author{Vincent Gerber, Tilman Hinnerichs}

% Name of the Course
\institute{Java Kurs}

% Fancy Logo
\titlegraphic{\hfill\includegraphics[height=1.25cm]{../templates/fsr_logo_cropped}}


\title{Java}
\subtitle{Collections part 2}
\date{\today}

\usepackage{color}

\begin{document}

\begin{frame}
\titlepage
\end{frame}

\begin{frame}{Overview}
\tableofcontents
\end{frame}

\section{Repetition}
\begin{frame}{Repetition}
	What we learned last time:
	\begin{itemize}
		\item How to use generics
		\item How to handle Javas lists, sets and iterators
	\end{itemize}
	What we will try to achieve today:
	\begin{itemize}
		\item How to use iterators on sets and lists
		\item How to use maps and what to with them
		\item What exceptions are and how to handle them
	\end{itemize}
\end{frame}

\begin{frame}{A quiz!}
	\begin{tabular}{c|c|c}
		&Set&List\\
		\hline
		\pause
		Same item twice in it? &&\\
		\hline
		Ordered?&&\\
		\hline
		Iterable?&&\\
		\hline
		What package to import &&\\
		\hline
		Declaring set type (variable type)&&\\
		\hline
		Building an instance (example)&&\\
		\hline
		Add an item&&\\
		\hline
		Removing an item&&\\
		
	\end{tabular}
\end{frame}

\begin{frame}{A quiz!}
	\begin{tabular}{c|c|c}
		&Set&List\\
		\hline
		Same item twice in it? &No!&Yes!\\
		\hline
		Ordered?&No!&Yes!\\
		\hline
		Iterable?&Yes!&Also yes!\\
		\hline
		What package to import &import java.util.*&import java.util.*\\
		\hline
		Declaring set type (variable type)&Set$<$T$>$ set&List$<$T$> $list\\
		\hline
		Building an instance (example)&= new HashSet$<$T$>$()& = new ArrayList$<$T$>$()\\
		\hline
		Add an item&set.add(item)&list.add(item)\\
		\hline
		Removing an item&set.remove(item)&list.remove(item)\\
	\end{tabular}
\end{frame}

\begin{frame}{Another quiz!}
	The iterator:\\
	\vspace{1cm}
	\begin{tabular}{l|l}
		&Iterator\\
		\hline
		How to declare&\\
		\hline
		How to build an instance&\\
		\hline
		First main function (With data type)&\\
		\hline
		Second main function (With data types)&\\
		\hline
		Third main function (With data type)&\\
		\hline
		How to get from collection?&\\
		
	\end{tabular}
\end{frame}

\begin{frame}[fragile]{How to iterate over sets and lists}
	\pause
	\begin{lstlisting}
		Set<T> mySet = new HashSet<T>();
		foreach(T item:mySet){
			item.doSomething();
		}
		
		List<T> myList = new ArrayList<T>();
		foreach(T item:myList){
		item.doSomething();
		}
	\end{lstlisting}
\end{frame}

\begin{frame}{Another quiz!}
	The iterator:\\
	\vspace{1cm}
	\begin{tabular}{l|l}
		&Iterator\\
		\hline
		How to declare&Iterator$<$T$>$ iter\\
		\hline
		How to build an instance&= new Iterator$<$T$>$()\\
		\hline
		First main function (With data type)&boolean iter.hasNext()\\
		\hline
		Second main function (With data types)&T iter.next()\\
		\hline
		Third main function (With data type)&T iter.remove()\\
		\hline
		How to get from collection(e.g. set)?&set.iterator()\\
		
	\end{tabular}
\end{frame}

\begin{frame}[fragile]{How to iterate over sets and lists using iterators}
	\begin{lstlisting}[basicstyle=\ttfamily\scriptsize]
	Set<T> mySet = new HashSet<T>();
	Iterator<T> myIter = mySet.iterator();
	
	while(myIter.hasNext()){
		T item = myIter.next();
		item.doSomething();
	}
	\end{lstlisting}
\end{frame}

\begin{frame}{Exercise}
	\begin{itemize}
		\item Create an array with 10 elements. Create a list and fill the list with the array 		elments. Create a set and fill the set with the list elments and create a map with the set elments as values and the index as key.
		\item Extend our vending machine with an internal storage
	\end{itemize}
\end{frame}

\section{Maps and iterators}
\begin{frame}[fragile]{Map}
	The interface \emph{Map} is not a subinterface of \emph{Collection}.\\
	A map contains pairs of key and value. Each key refers to a value. 
	Two keys can refer to the same value. There are not two equal keys in one map.
	\emph{Map} is part of the package \texttt{java.util}.
	\vfill
	\begin{lstlisting}[basicstyle=\ttfamily\scriptsize]
	public static void main (String[] args) {
	
	Map<Integer, String> map = 
	new HashMap<Integer, String>();
	
	map.put(23, "foo");
	map.put(28, "foo");
	map.put(31, "bar");
	map.put(23, "bar"); // "bar" replaces "foo" for key = 23
	
	System.out.println(map);
	// prints: {23=bar, 28=foo, 31=bar}
	}
	\end{lstlisting}
\end{frame}

\begin{frame}[fragile]{Key, Set and Values}
	You can get the set of keys from the map.
	Because one value can exist multiple times a collection is used for the values.
	\begin{lstlisting}[basicstyle=\ttfamily\scriptsize]
	public static void main (String[] args) {
	
	// [...] map like previous slide
	
	Set<Integer> keys = map.keySet();
	Collection<String> values = map.values();
	
	System.out.println(keys);
	// prints: [23, 28, 31]
	
	System.out.println(values);
	// prints: [bar, foo, bar]
	}
	\end{lstlisting}
\end{frame}

\begin{frame}[fragile]{Iterator}
	To iterate over a map use the iterator from the set of keys.
	\begin{lstlisting}[basicstyle=\ttfamily\scriptsize]
	public static void main (String[] args) {
	
	// [...] map, keys, values like previous slide    	    
	Iterator<Integer> iter = keys.iterator();
	
	while(iter.hasNext()) {
	System.out.print(map.get(iter.next()) + " ");
	} // prints: bar foo bar
	
	System.out.println(); // print a line break
	
	for(Integer i: keys) {
	System.out.print(map.get(i) + " ");
	} // prints: bar foo bar
	}
	\end{lstlisting}
\end{frame}

\begin{frame}[fragile]{Nested Maps}
	Nested maps offer storage with key pairs.
	\begin{lstlisting}[basicstyle=\ttfamily\scriptsize]
	public static void main (String[] args) {		
	
	Map<String, Map<Integer, String>> addresses = 
	new HashMap<String, Map<Integer, String>>();
	
	addresses.put("Noethnitzer Str.", 
	new HashMap<Integer, String>());
	
	addresses.get("Noethnitzer Str.").
	put(46, "Andreas-Pfitzmann-Bau");
	addresses.get("Noethnitzer Str.").
	put(44, "Fraunhofer IWU");
	}
	\end{lstlisting}
\end{frame}

\begin{frame}[fragile]{Maps and For Each}
	You can interate through the entry set of a map (available before Java 1.8)
	\begin{lstlisting}[basicstyle=\ttfamily\scriptsize]
	Map<String, String> map = ...
	for (Map.Entry<String, String> entry : map.entrySet()) {
	System.out.println("Key: " + entry.getKey() +
	", value" + entry.getValue());
	}
	\end{lstlisting}
\end{frame}

\begin{frame}{Overview}
	\begin{center}
		\begin{tabular}{ l | l }
			List & $ \bullet $ Keeps order of objects \\
			& $ \bullet $ Easily traversible \\
			& $ \bullet $ Search not effective \\
			\hline
			Set  & $ \bullet $ No duplicates \\
			& $ \bullet $ No order - still traversible \\
			& $ \bullet $ Effective searching \\
			\hline
			Map  & $ \bullet $ Key-Value storage \\
			& $ \bullet $ Search super-effective \\
			& $ \bullet $ Traversing difficult
			
		\end{tabular}
	\end{center}
\end{frame}

\begin{frame}{Easy and some more complex exercises}
	\huge Hier könnten Ihre Aufgaben zu Mpas und Iteratoren stehen!
\end{frame}


\end{document}
